%-------------------------
% Cover Letter - Cerebras New Grad (Inference Frontend)
% Author: Ajay Venkatesh
%-------------------------

\documentclass[letterpaper,11pt]{article}

\usepackage[top=0.8in, bottom=0.8in, left=0.9in, right=0.9in]{geometry}
\usepackage[hidelinks]{hyperref}
\usepackage{lmodern}
\usepackage{parskip}
\usepackage{microtype}

\setlength{\parskip}{6pt}
\setlength{\parindent}{0pt}

\begin{document}

\begin{flushleft}
\textbf{\large Ajay Venkatesh} \\
Brooklyn, NY \\
+1 516 585 9013 \\
\href{mailto:av3855@nyu.edu}{av3855@nyu.edu}
\end{flushleft}

\vspace{10pt}

May 4, 2026

\vspace{6pt}

Hiring Manager \\
Cerebras Systems

\vspace{10pt}

Dear Hiring Manager,

My name is Ajay Venkatesh --- finishing my M.S. in Computer Engineering at NYU (3.9 GPA, May 2026) with roughly four years of production software experience: three at PwC in Bengaluru, a summer SDE internship at Amazon, and ongoing on-campus work at NYU's Novel AI Technologies. My graduate coursework has concentrated at the systems-and-ML intersection --- \textbf{Computer Architecture, Machine Learning, Deep Learning, Big Data, and System Design} --- which is the same layer-cake the Cerebras new-grad posting calls for.

The role's framing --- working across hardware, compilers, distributed systems, and ML frameworks --- maps to where I've been pushing my own learning. In my Big Data course I built a \textbf{PySpark ML pipeline on HDFS} over \textbf{7M flight records}, training Random Forest and Gradient Boosted Tree models to \textbf{85\%+} accuracy; small scale next to a wafer-scale engine, but the same shape of problem: feature pipelines, distributed compute, and squeezing accuracy out of irregular data. Deep Learning coursework gave me the model-side intuition (CNNs, transformers, optimizer behavior); Computer Architecture filled in what's actually happening underneath the abstractions.

On the production side, my \textbf{C++} foundation is from coursework rather than industry, but I've shipped systems work in adjacent stacks --- a distributed event-driven backend on \textbf{Azure} at PwC over SQL Server, \textbf{serverless infrastructure on AWS CDK} at Amazon (API Gateway, Lambda, DynamoDB, multi-region CI/CD), and parallelized \textbf{ETL pipelines} at NYU that cut processing time \textbf{60\%}. I'm most useful where backend systems meet ML, and least afraid of dropping a layer down when something doesn't behave the way the abstraction said it would.

What pulls me toward Cerebras is the conviction behind the WSE --- that inference throughput and latency are bottlenecked by general-purpose hardware, and the answer is a stack rebuilt from silicon up. I'd like to learn how that stack actually works, end-to-end, alongside the people building it.

Thanks for considering my application --- I'd welcome the chance to discuss further.

\vspace{12pt}

Sincerely, \\
Ajay Venkatesh

\end{document}
